% =============================================================================
% PASTE-READY THESIS SNIPPET — generative-model benchmark
% Draft only. The thesis itself is under a freeze policy (edit on request only),
% so this lives in docs/ as a snippet to copy in when you decide where it goes.
%
% Figure: results/benchmark_baselines/fig_benchmark_publication.pdf
% Numbers verified against results/benchmark_baselines/rmse_table.json (2026-06-11)
% Generated by scripts/analysis/benchmark_{baselines,diffusion,flow_matching}.py
%            scripts/figures/plot_benchmark_publication.py
% =============================================================================

% ---- Figure environment -----------------------------------------------------
\begin{figure}[t]
  \centering
  \includegraphics[width=\textwidth]{figures/fig_benchmark_publication.pdf}
  \caption[Predictive benchmark of the gLV model against ML and generative
  baselines]{\textbf{The mechanistic gLV model outperforms machine-learning and
  generative baselines on leave-one-patient-out prediction.}
  \textbf{(a)}~LOO-CV RMSE (mean~$\pm$~SD over the ten Dieckow patients) for the
  one-step-ahead guild-composition prediction task, grouped by model family:
  mechanistic ODE (gLV variants), ML regression (Random Forest, Ridge, LASSO,
  and the persistence / train-mean references), and generative models
  (conditional Flow Matching and DDPM). \textbf{(b)}~Paired per-patient
  comparison of the best model (gLV+AGORA) against one representative of each
  competing family; each grey line connects the two RMSE values for one held-out
  patient. gLV wins on all ten patients against every competitor
  (one-sided Wilcoxon signed-rank $p=0.001$ in each case, the minimum attainable
  at $n=10$).}
  \label{fig:generative-benchmark}
\end{figure}

% ---- Body text (place near the LOO-CV / model-validation section) -----------
To test whether the mechanistic structure of the gLV model is actually
necessary—or whether a purely data-driven predictor would do as well—we
benchmarked it against two families of baselines on the same
leave-one-patient-out (LOO) one-step-ahead prediction task. The first family is
classical machine-learning regression (Random Forest, Ridge, LASSO), together
with two naive references (persistence, $\bm{\phi}(t{+}1)=\bm{\phi}(t)$, and the
training-set mean composition). The second family is modern generative models:
a conditional denoising diffusion model (DDPM) and conditional Flow Matching,
the ODE-based successor to diffusion. All methods were trained and evaluated
under an identical protocol—nine patients for training ($18$ week-to-week
transitions), the tenth held out, with the third week predicted by chaining the
model from the predicted second week.

The mechanistic model dominates (Fig.~\ref{fig:generative-benchmark}a). gLV+AGORA
attains an LOO-CV RMSE of $0.050\pm0.017$, roughly $1.8\times$ lower than the
best ML baseline (Random Forest, $0.091\pm0.022$) and $2$--$3\times$ lower than
the generative models (Flow Matching $0.107\pm0.025$; DDPM $0.166\pm0.031$). The
advantage is consistent at the level of individual patients
(Fig.~\ref{fig:generative-benchmark}b): gLV+AGORA achieves the lowest RMSE on all
ten held-out patients against every competitor, a result that is significant at
the floor of the one-sided Wilcoxon signed-rank test for $n=10$ ($p=0.001$).

Two secondary observations support the interpretation that the data regime, not
the model class, is decisive. First, Flow Matching outperforms DDPM by $36\%$,
mirroring the training-stability advantage of the ODE (probability-flow)
formulation over the stochastic reverse process---yet both still lose to gLV.
Second, augmenting the ML training set with synthetic trajectories sampled from a
diffusion model \emph{degraded} Ridge and Random Forest rather than improving
them, indicating that with only eighteen real transitions a generative model
cannot recover the underlying dynamics and instead injects noise. Together these
results show that, in the small-sample longitudinal regime characteristic of
clinical oral-microbiome studies, encoding the ecological interaction structure
explicitly is worth far more than additional model flexibility---the central
motivation for the mechanistic approach taken in this thesis.
